% Shared preamble for Quiz 04 (quiz04.tex and solutions.tex).
% Compile with: xelatex -interaction=nonstopmode <file>.tex   (run twice)
%
% Same house look as the pen-and-paper sheets in
% lecture-note/*/pen-and-paper/: Charter, A4, generous leading, room to draw.

\usepackage[top=0.7in, bottom=0.6in, left=0.8in, right=0.8in]{geometry}
\usepackage{amsmath}
\usepackage{amssymb}
\usepackage{array}
\usepackage{graphicx}
\usepackage[hidelinks]{hyperref}
\usepackage{tcolorbox}
\usepackage{enumitem}
\usepackage{etoolbox}
\usepackage{tikz}
\usepackage{fontspec}
\usetikzlibrary{calc,positioning,arrows.meta}
\definecolor{pltblue}{HTML}{1F77B4}
\definecolor{rulegrey}{HTML}{6A6D75}

% Body face. Charter, and the four ways to get it. XCharter *is* Charter and
% ships with a full TeX Live, so it is asked for first and by file: a CI runner
% answers yes to \IfFontExistsTF{Charter}, then cannot typeset it.
\IfFontExistsTF{XCharter-Roman.otf}{
  \setmainfont{XCharter}[
    Extension      = .otf,
    UprightFont    = *-Roman,
    BoldFont       = *-Bold,
    ItalicFont     = *-Italic,
    BoldItalicFont = *-BoldItalic]
}{
  \IfFontExistsTF{Charter}{\setmainfont{Charter}}{
    \IfFontExistsTF{Iowan Old Style}{\setmainfont{Iowan Old Style}}{
      \IfFontExistsTF{Georgia}{\setmainfont{Georgia}}{}}}}

\setlength{\parindent}{0pt}
\setlength{\parskip}{0.4em}
\setlength{\emergencystretch}{1.2em}
\raggedbottom
\pagestyle{empty}

% ---------------------------------------------------------------------------
% The submission link.
%
% Both the QR square and the printed address are the course's own short link,
% not the form's: paper outlives any one form. Caddy on the course droplet
% points go.skojaku.com/ans-quiz04 at the current Google Form
% (/etc/caddy/Caddyfile, block go.skojaku.com). Swapping the form is a line on
% the server, not a reprint.
%
% If the link ever changes, rebuild the square with
%
%   uvx --from segno segno --output=quiz04-form-qr.png --scale=20 --border=1 \
%       --error=m "<the url>"
%
% PNG rather than PDF for the square: segno's PDF is too spare for xdvipdfmx to
% include, and a QR is squares, so 700px across a 3cm box is already sharper
% than the printer.
% ---------------------------------------------------------------------------
\newcommand{\formlink}{https://go.skojaku.com/ans-quiz04}
\newcommand{\formshort}{\href{\formlink}{\texttt{go.skojaku.com/ans-quiz04}}}

% Quiz header. #1 = title, #2 = the line under it, #3 = term.
\newcommand{\quizhead}[3]{%
  {\footnotesize\textsc{Advanced Network Science} \hfill \textsc{#3}}\par
  \vspace{-0.2em}\rule{\textwidth}{1pt}\par
  \vspace{0.3em}
  {\LARGE\bfseries #1}\par
  \vspace{0.1em}
  {\small #2}\par
  \vspace{0.7em}
  {\small Name \underline{\hspace{6cm}} \hfill Date \underline{\hspace{3.2cm}}}\par
  \vspace{0.7em}}

% A question heading with its point value.
\newcommand{\question}[3]{%
  \vspace{0.5em}
  {\large\bfseries Question #1 \quad\normalsize\mdseries (#2 points)}\par
  \vspace{-0.3em}{\color{rulegrey}\rule{\textwidth}{0.6pt}}\par
  \vspace{0.1em}
  {\bfseries #3}\par
  \vspace{0.2em}}

% An empty box to draw or write in. #1 = height, #2 = caption under the rule.
\newcommand{\drawbox}[2]{%
  \begin{tcolorbox}[colback=white, colframe=rulegrey, boxrule=0.5pt,
                    sharp corners, height=#1, valign=top,
                    boxsep=2pt, left=4pt, right=4pt, top=3pt, bottom=3pt]
  {\footnotesize\color{rulegrey} #2}
  \end{tcolorbox}}

% Ruled writing lines, #1 of them.
\newcommand{\writelines}[1]{%
  \vspace{0.3em}
  \foreach \i in {1,...,#1}{%
    {\color{rulegrey}\rule{\textwidth}{0.4pt}}\par\vspace{1.5em}}}

% ---------------------------------------------------------------------------
% Log--log graph paper for the CCDF. x runs over three decades (k = 1 to
% 1000), y over four (10^-4 to 1). One decade is the same length on both
% axes, so a slope measured with a ruler in cm is the slope in decades.
% #1 = the grid's letter, #2 = anything drawn on it (the key puts the answer
% here), in axis units: x = log10 k, y = log10 of the fraction.
% ---------------------------------------------------------------------------
\tikzset{
  majorgrid/.style = {line width=0.55pt, black!45},
  minorgrid/.style = {line width=0.35pt, black!18},
  loglab/.style    = {font=\footnotesize},
  logminor/.style  = {font=\tiny, black!55},
}
\newcommand{\loggrid}[2]{%
\begin{tikzpicture}[x=1.95cm, y=1.95cm]
  \foreach \d in {0,1,2} \foreach \m in {2,...,9}
    \draw[minorgrid] ({\d+log10(\m)},-4) -- ({\d+log10(\m)},0);
  \foreach \d in {-4,-3,-2,-1} \foreach \m in {2,...,9}
    \draw[minorgrid] (0,{\d+log10(\m)}) -- (3,{\d+log10(\m)});
  \foreach \i in {0,1,2,3} \draw[majorgrid] (\i,-4) -- (\i,0);
  \foreach \j in {-4,...,0} \draw[majorgrid] (0,\j) -- (3,\j);
  \foreach \i/\lab in {0/1, 1/10, 2/100, 3/1000}
    \node[loglab, below] at (\i,-4) {\lab};
  \foreach \d/\a/\b in {0/2/5, 1/20/50, 2/200/{}} {
    \node[logminor, below] at ({\d+log10(2)},-4) {\a};
    \node[logminor, below] at ({\d+log10(5)},-4) {\b};}
  \foreach \j/\lab in {-4/{0.0001}, -3/{0.001}, -2/{0.01}, -1/{0.1}, 0/{1}}
    \node[loglab, left] at (0,\j) {\lab};
  \foreach \d in {-4,-3,-2,-1} {
    \node[logminor, left] at (0,{\d+log10(2)}) {2};
    \node[logminor, left] at (0,{\d+log10(5)}) {5};}
  \node[loglab] at (1.5,-4.45) {degree $k$};
  \node[loglab, rotate=90] at (-0.66,-2) {fraction with degree $> k$};
  \node[font=\bfseries, anchor=south west] at (0,0.02) {#1};
  #2
\end{tikzpicture}}
